\chapter{Étude de l'existant et positionnement}
\label{ch:etat-art}

\section{Périmètre de l'étude}

Cette étude compare LexHub aux offres accessibles à un cabinet marocain, selon une grille retenue au premier suivi PFA : atelier voix~$\rightarrow$~plaidoirie, gestion de cabinet, SaaS multi-tenant, travail en équipe.

Sources consultées : \texttt{smartavocat.ma}, \texttt{avocato.ma}, \texttt{avocatpro.ma}, \texttt{legaliq.ma}, \texttt{kanoun.com}, \texttt{adala.ai}, \texttt{diskod.ma}.

\section{Analyse comparative (benchmarking)}
\label{sec:benchmark}

Le tableau~\ref{tab:benchmark} indique, pour chaque offre, ce qu'elle couvre et ce qui manque par rapport au besoin LexHub : une plaidoirie vocale intégrée au cabinet.

\begin{table}[H]
\centering
\caption{Benchmarking marché --- applications existantes vs besoin LexHub}
\label{tab:benchmark}
\footnotesize
\renewcommand{\arraystretch}{1.35}
\begin{tabular}{@{}p{2.4cm}p{4.6cm}p{5.6cm}@{}}
\toprule
\textbf{Offre} & \textbf{Services offerts} & \textbf{Ce qui manque face à LexHub} \\
\midrule
SmartAvocat \newline {\scriptsize smartavocat.ma} &
Gestion cabinet, dossiers bilingues &
Pas de flux plaidoirie vocale \\
\midrule
Avocato \newline {\scriptsize avocato.ma} &
Gestion + audiences / Mahakim &
Pas de parcours dictée $\rightarrow$ document métier \\
\midrule
AvocatPro \newline {\scriptsize avocatpro.ma} &
Gestion + Mahakim, organisation &
Pas de cœur plaidoirie \\
\midrule
Legaliq \newline {\scriptsize legaliq.ma} &
IA juridique / fiscale, documents &
Pas un SaaS cabinet complet \\
\midrule
Kanoun \newline {\scriptsize kanoun.com} &
IA droit marocain + OHADA &
Pas multi-tenant cabinet \\
\midrule
Adala \newline {\scriptsize adala.ai} &
Assistant IA, Q\&A, documents &
Pas d'équipe cabinet \\
\midrule
Dragon Legal / Diskod \newline {\scriptsize diskod.ma} &
Dictée desktop (Nuance) &
Pas SaaS, pas dossiers LexHub \\
\midrule
\textbf{LexHub} \newline {\scriptsize lexhub.space} &
Atelier 5 étapes ; dossiers ; équipe ; SaaS &
--- \\
\bottomrule
\end{tabular}
\end{table}

\section{Positionnement de la solution LexHub}

LexHub, tel qu'il est en ligne, combine un SaaS de cabinet (isolation, rôles, dossiers), un atelier \textbf{Dossier \& pièces / Dictée / Transcription / Rédaction IA / Export}, et un document Word / PDF. L'IA assiste ; elle n'est pas un chatbot autonome.

\section{Synthèse et choix retenus}

\begin{enumerate}[leftmargin=1.4em]
  \item Un outil de cabinet, pas un chatbot générique.
  \item Un cœur produit en cinq étapes voix~$\rightarrow$~document, distinct d'un logiciel Mahakim.
  \item La fidélité à la dictée : aucun fait, article ou montant inventé.
  \item Un SaaS facturable (essai, plans, packs, sièges), détaillé au chapitre~\ref{ch:economie}.
\end{enumerate}
